\documentclass[11pt]{article}
\pagestyle{empty}

%\setlength{\topmargin}{.25in}

\setlength{\textwidth}{6in}
\setlength{\oddsidemargin}{.25in}
\setlength{\topmargin}{0in}
\setlength{\textheight}{8.25in}


\renewcommand{\baselinestretch}{1.1}
\usepackage{amsmath}
\usepackage{epsfig}
\usepackage{amssymb}
\usepackage{amscd}
\usepackage{times}
\usepackage{graphics}
%\usepackage{chicago}
%\usepackage{fullpage}
%\usepackage{apalike}
\newtheorem{proposition}{Proposition}
\newtheorem{corollary}{Corollary}
\newtheorem{theorem}{Theorem}
\newtheorem{lemma}{Lemma}
\newtheorem{example}{Example}
\newtheorem{claim}{Claim}
\newtheorem{assumption}{Assumption}
\newtheorem{definition}{Definition}
\def\argmax{\mathop{\rm argmax}\limits}
\def\max{\mathop{\rm max}\limits}


\begin{document}

\noindent \large Midterm 2017, ``The Mathematical Foundations of Political Methodology"

\normalsize
\ \\You have 90 minutes for this exam.  You may use a calculator and refer to your Pemberton \& Rau book.  Please do not use the internet during the exam.  Email your solutions, \textit{with your work shown}, to Milena.  Solutions without work will get no credit.

\begin{enumerate}

\item Given the following hold, what do we know about the relationship between sets $A$ and $C$?

\begin{enumerate}
\item $A\cap B= B$ and $B\cap C\not=\emptyset$


\ \\ANSWER:  $A\cap C\not=\emptyset$.

\item $A\cup B= B$ and $B\cap C\not=\emptyset$


\ \\ANSWER:  Nothing.

\item $B \cap D=C$ and $B\cap C=A$

\ \\ANSWER:   $A = C$.

\end{enumerate}

\vspace{.2in}

\item Find the limit of each of the following sequences as $n\rightarrow \infty$, if the limit exists.
\begin{enumerate}
\item $${{4n^2-1}\over{(n+1)^2}}$$

\ \\ANSWER:  It's 4.

\item $${{(-1)^n n^{1\over 2}}\over{7+n^{3\over 2}}}$$


\ \\ANSWER:   It's 0.

\item $${{(-1)^n n^{1\over 2}}\over{n^{1\over 4}+n^{1\over 2}}}$$

\ \\ANSWER:   It doesn't exist.
\end{enumerate}

\vspace{.2in}

\item Use the $(\delta, \epsilon)$ definition of a limit to prove that the function $$f(x)={2\over 3}x-17$$ is everywhere continuous.

\ \\ANSWER:  We need that the limit as $x\rightarrow x_o$ of $f(x)$ equals ${2\over 3}x_o-17$ for all $x_o$.  For any $\epsilon>0$, we can find a $\delta>0$ such that $|{2\over 3}x-17-{2\over 3}x_o+17|<\epsilon$ whenever $|x-x_o|<\delta$.  Rearranging, we get ${2\over 3}|x-x_o|<\epsilon$ when $|x-x_o|<\delta$.  Setting $\delta={3\over 2} \epsilon$ proves the result.

\vspace{.2in}
\item Let $$f(x)={{x}\over{x^2+1}}+5$$
\begin{enumerate}
\item Find the critical points of $f$ and characterize them as maxima, minima, or inflection points.

\ \\ANSWER:  Critical points at $\pm 1$; 1 is a local maximum, $-1$ is a local minimum.

\item Find the regions on which $f$ is convex and concave.

\ \\ANSWER:  $f$ is concave on $(-\infty, -\sqrt{3})$, convex on $(-\sqrt{3}, 0)$, concave on $(0, \sqrt{3})$, and convex on $(\sqrt{3}, \infty)$.  We can find this by evaluating the second derivative. 

\item Find any asymptotes that $f$ may have.

\ \\ANSWER:   As $x\rightarrow \infty$ $f$ approaches the line $y=5$ (from above). As $x\rightarrow -\infty$ $f$ approaches the line $y=5$ (from below). 

\end{enumerate}

\item Perform the following matrix operations:

\begin{enumerate}

\item Find the inverse of $$\begin{bmatrix}1 & 1& 2\\ 0& 1& 1\\0&0&1
\end{bmatrix}$$


\ \\ANSWER:   It is $$\begin{bmatrix}1 & -1& -1\\ 0& 1& -1\\0&0&1 \end{bmatrix}$$

\item Find the rank of 
$$\begin{bmatrix} 
1& 2& 5\\
-3& 2& -1\\
0& 4& 7\\
2& 0& 3
\end{bmatrix}$$

\ \\ANSWER:  It is 2.

\item Check whether the following matrix is positive definite, positive semidefinite, negative definite, negative semidefinite, or none of the above.
$$\begin{bmatrix}
1& 3& 2\\3& 2& 1\\2& 1& 1
\end{bmatrix}$$

\ \\ANSWER:  It is none of the above (the 2X2 principal minors are different signs).

\end{enumerate}

\vspace{.2in}

\item Find and classify all of the critical points of the following functions:

\begin{enumerate}

\item  $$f(x, y)=-10x^2+4xy-{{2y^2}\over 5}$$


\ \\ANSWER:  There are a continuum of critical points: all points of the form $(C, 5C)$ with $C\in \mathbb{R}$.  They are all global optima, because $f$ is globally concave: Its Hessian is $$\begin{bmatrix}-20 & 4\\4& -{4\over 5}\end{bmatrix}$$ so it has determinant 0 and its diagonals are negative.  Thus the Hessian is negative semidefinite.

\item $$f(x, y, z)=-x^2+{\log}(x)-{1\over y}-y-z^2+z \hspace{.1in}\hbox{ for }x, y, z>0$$

\ \\ANSWER:   There is one critical point: $({1\over \sqrt{2}}, 1, {1\over 2})$.  It is a global optimum because the function is the sum of concave functions of one variable.  No Hessian required.

\item $$f(x, y)=x^2+y^2-xy \hbox{ subject to the constraint that }2x-y=3$$

\ \\ANSWER:  The Lagrangian is $$L(x, y, \lambda)=x^2+y^2-xy-\lambda(2x-y-3).$$  The first order conditions imply that $2x-y-2\lambda=0$, $2y-x+\lambda=0$, and $2x-y-3=0$.  Solving for $x$ and $y$ we get $({3\over 2},0)$ as a critical point (incidentally, $\lambda={3\over 2}$).  It is a constrained (global) minimum.  We know this because $x^2+y^2-xy$ is globally convex and our constraint is linear. The Hessian of $f$ is $$\begin{bmatrix}2& -1\\-1& 2\end{bmatrix}$$ which has determinant $3>0$ and positive diagonals. 
\end{enumerate}

\end{enumerate}

\end{document}
